\section{Related Work}
\label{sec:related}

\paragraph{EEG foundation models.}
Three open-weight EEG foundation models anchor our evaluation.
\textbf{LaBraM} \citep{labram2024} is a ViT with
channel-position embeddings over the 10--20 system and a vector-quantised
tokenizer, pretrained with masked patch prediction on roughly 2{,}000
hours of mixed EEG. \textbf{CBraMod} \citep{cbramod2025} replaces
single-stream attention with a criss-cross transformer that factorises
spatial and temporal attention, and uses a CNN-plus-FFT spectral patch
embedding; the encoder is pretrained on TUEG with masked patch
modelling. \textbf{REVE} \citep{reve2025} is a ViT with a 4D Fourier
positional encoding pretrained by masked autoencoding on $25{,}000$
subjects across heterogeneous montages. The three backbones span
complementary design choices --- tokenization strategy (VQ vs.\
continuous), attention factorisation, and pretraining scale --- and
have each been reported to achieve strong downstream performance on
event-related EEG \citep{eegfmbench2025,adabrain2025}. They also have
materially different input assumptions: LaBraM expects z-scored
inputs, whereas CBraMod and REVE are calibrated to raw microvolt-scale
signals. This per-model input convention is neither documented in
aggregate benchmarks nor invariant to preprocessing, and we will show
in Sec.~\ref{sec:results} that it interacts with causal anchor
ablation in a way that would be invisible under a single global
normalisation.

\paragraph{Subject leakage and evaluation protocols.}
A distinct line of work has audited the evaluation protocols under
which EEG deep-learning results are reported. Large-scale audits have
established that segment-based or trial-level cross-validation in
clinical EEG inflates accuracy by exposing models to segments from the
same subject at train and test time
\citep{brookshire2024leakage,smalln2025biorxiv}. Subject-disjoint
cross-validation is increasingly adopted as the primary protocol, and
benchmark efforts such as EEG-FM-Bench \citep{eegfmbench2025},
AdaBrain-Bench \citep{adabrain2025}, EEG-Bench
\citep{eegbench2025}, and Brain4FMs \citep{brain4fms2026} enforce some
form of cross-subject splitting. These audits establish \emph{that}
subject leakage inflates reported numbers and document the size of the
inflation; they do not predict, from dataset metadata alone, which
datasets will collapse and which will retain a signal once leakage is
removed. Our $2\times 2$ factorial takes the subject-disjoint
protocol as the starting point and asks a complementary
question --- under an already-honest protocol, which structural
properties of a dataset (subject relation of the label, linear-probe
separability) predict whether FM features carry a usable signal at
all.

\paragraph{Neural anchors for clinical rsEEG labels.}
A parallel neurophysiological literature asks which spectral and
band-specific features carry trait- or state-level clinical signal
in resting-state EEG. Classical oscillatory accounts link theta and
alpha rhythms to cognitive and attentional state
\citep{klimesch1999,klimesch2012update,pfurtscheller1979}, and more
recent work has quantified the discriminative power of band-ratio
features for cognitive stress, typically finding that a handful of
scalar ratios (theta/alpha ratio, frontal alpha asymmetry) carry most
of the classifiable variance in cross-subject designs. The neural
efficiency and allostatic regulation literatures further motivate a
ceiling argument for trait-like resting-state paradigms in young,
high-cognitive-reserve cohorts \citep{lovibond1995}. These works
characterise \emph{signal-side} anchors --- which band or ratio
discriminates --- but do not ask the downstream question of whether
a pretrained FM's representation has actually learned that anchor.
Our causal anchor ablation (Sec.~\ref{sec:results}) operationalises
this check by intervening on the FOOOF-parameterised aperiodic
$1/f$ background \citep{donoghue2020fooof} and on narrow-band
power independently, and reporting which ablation collapses each
FM's probe.

\paragraph{Aperiodic $1/f$ activity as subject identity signal.}
A recent strand of work has drawn attention to the aperiodic $1/f$
component of the EEG power spectrum, parameterised separately from
periodic peaks by FOOOF \citep{donoghue2020fooof}. Aperiodic slope
and offset are stable within an individual across recording sessions
and carry idiosyncratic information that distinguishes subjects
independently of task-evoked oscillations, paralleling the broader
observation that functional-connectivity patterns fingerprint
individuals \citep{finn2015,campisi2014,marcel2007}. At the
representation level, SSL-pretrained EEG backbones have been shown to
encode strong subject identity --- a property that has been framed as
useful signal for within-subject BCI calibration, where
intra-subject regularity is the target \citep{multidataset2025},
but which behaves as a confound for cross-subject clinical
screening. Our work treats this asymmetry explicitly: we use
FOOOF-based aperiodic ablation not as a preprocessing step but as an
intervention whose purpose is to ask whether each FM's downstream
probe is anchored on the $1/f$ axis, on periodic peaks, on broadband
combinations of the two, or on no recoverable spectral anchor.

\paragraph{Factorial and taxonomy-style analyses in EEG benchmarking.}
Existing aggregate benchmarks organise EEG-FM evaluation either as a
single-dataset case study or as a leaderboard that aggregates across
many tasks and datasets. EEG-FM-Bench \citep{eegfmbench2025},
AdaBrain-Bench \citep{adabrain2025}, EEG-Bench \citep{eegbench2025},
and Brain4FMs \citep{brain4fms2026} each report dozens of
task--dataset cells and broadly converge on the observation that
FMs are not uniformly superior across clinical cross-subject tasks
\citep{kauffmann2025critique,eegchallenge2025}. Their taxonomies,
however, are organised along axes largely orthogonal to ours: task
category (motor imagery, SSVEP, affect, disease), pretraining
objective (masked reconstruction, contrastive, diffusion), or
fine-tuning strategy (frozen, LoRA, full-parameter). None separates
datasets by \emph{subject relation of the label} (within-subject
paired vs.\ between-subject label) crossed with \emph{linear-probe
separability under subject-disjoint CV} (separable vs.\
null-indistinguishable), and none asks whether the quadrants of this
factorial exhibit distinguishable mechanisms. Our work is not a
bigger benchmark --- four datasets and three backbones is deliberately
small --- but a structural re-organisation that foregrounds
\emph{why} benchmark numbers fragment across cohorts of comparable
size, paired with a diagnostic toolkit whose scope conditions state
in advance which cells each tool is valid in.
